\begin{alist}
\item Από την γνωστή ανισότητα $-1 \le \eta\mu x \le 1$ προκύπτει ότι $-2 \le 2\eta\mu x \le 2$, οπότε $-3 \le 2\eta\mu x - 1 \le 1$, δηλαδή $-3 \le f(x) \le 1$. Επιπλέον στο διάστημα $[0, 2\pi]$ η ισότητα $f(x) = -3$ ισχύει όταν $\eta\mu x = -1$ δηλαδή $x = \dfrac{3\pi}{2}$, ενώ η $f(x) = 1$ ισχύει όταν $\eta\mu x = 1$ δηλαδή $x = \dfrac{\pi}{2}$. \\
Επομένως η $f$ παρουσιάζει:
\begin{itemize}
\item Ολικό ελάχιστο για $x = \dfrac{3\pi}{2}$, το $f\left(\dfrac{3\pi}{2}\right) = -3$.
\item Ολικό μέγιστο για $x = \dfrac{\pi}{2}$, το $f\left(\dfrac{\pi}{2}\right) = 1$.
\end{itemize}

\item Με $x=0$ έχουμε $f(0) = -1$, οπότε η $\text{C}_f$ τέμνει τον άξονα $y'y$ στο σημείο $\text{A}(0, -1)$. \\
Με $y=f(x)=0$ έχουμε $2\eta\mu x - 1 = 0 \Leftrightarrow \eta\mu x = \dfrac{1}{2}$, οπότε για $x \in [0, 2\pi]$ βρίσκουμε $x = \dfrac{\pi}{6}$ ή $x = \dfrac{5\pi}{6}$. \\
Επομένως τα κοινά σημεία της $\text{C}_f$ με τον άξονα $x'x$ είναι τα $\text{B}\left(\dfrac{\pi}{6}, 0\right), \Gamma\left(\dfrac{5\pi}{6}, 0\right)$.

\item Συμπληρώνουμε τον παρακάτω πίνακα τιμών.

\begin{center}
\begin{tabular}{|l|l|l|l|l|l|}
\hline
$x$ & $0$ & $\dfrac{\pi}{2}$ & $\pi$ & $\dfrac{3\pi}{2}$ & $2\pi$ \\[8pt] \hline
$\eta\mu x$ & $0$ & $1$ & $0$ & $-1$ & $0$ \\[3pt] \hline
$f(x)$ & $-1$ & $1$ & $-1$ & $-3$ & $-1$ \\[3pt] \hline
\end{tabular}
\end{center}

Με τη βοήθεια του παραπάνω πίνακα προκύπτει η επόμενη γραφική παράσταση $\text{C}_f$ της $f$.

\begin{center}
\begin{tikzpicture}
\begin{axis}[
    axis lines=middle,
    xmin=-0.5, xmax=6.5,
    ymin=-3.5, ymax=2.5,
    xtick={1.57, 3.14, 4.71, 6.28},
    xticklabels={$\pi/2$, $\pi$, $3\pi/2$, $2\pi$},
    ytick={-3, -2, -1, 1, 2},
    grid=both,
    grid style={dashed, gray!40},
    thick,
    width=12cm, height=8cm,
    samples=100
]
\addplot[domain=0:2*pi, maincolor, very thick] {2*sin(deg(x)) - 1} node[pos=0.15, above, text=black] {$\text{C}_f$};
\addplot[mark=*, fill=black] coordinates {(0, -1)};
\addplot[mark=*, fill=black] coordinates {(2*pi, -1)};
\end{axis}
\end{tikzpicture}
\end{center}

\item Από την ισότητα $f(\alpha) = f\left(\dfrac{\pi}{2} - \alpha\right)$, έχουμε:
\end{alist}
